\Subhead{1.10.1}{Распределённые автоколебательные системы. Лазер. Условия самовозбуждения}
\textbf{Физическая схема.} Активная среда с инверсией населённостей усиливает поле посредством вынужденного излучения. Резонатор обеспечивает обратную связь и пространственные моды, накачка поддерживает инверсию, потери и насыщение определяют порог и стационарную мощность. Это распределённая система: поле и поляризация зависят от координат, а модальная модель является приближённым сокращением полевых уравнений.

\paragraph{Условие самовозбуждения за обход.}
Рассмотрим линейный резонатор длины $L$ с зеркалами интенсивностного отражения $R_1,R_2$, активным участком длины $l_g$. Пусть $g$ --- коэффициент усиления \emph{интенсивности} на единицу длины в активном участке, $\alpha_{\mathrm{int}}$ --- распределённые потери интенсивности по всему резонатору. После полного обхода интенсивность умножается на
\begin{equation}
 \mathcal G=R_1R_2\exp(2gl_g-2\alpha_{\mathrm{int}}L).
\end{equation}
При малом поле генерация начинается, когда усиление превышает потери, а фаза обхода равна $2\pi q$. Порог
\begin{equation}\label{eq:laserthreshold}
 g_{\mathrm{th}}=\frac{\alpha_{\mathrm{int}}L+\tfrac12\ln[(R_1R_2)^{-1}]}{l_g}.
\end{equation}
Формула предполагает отсутствие дополнительных локальных потерь; их логарифмические вклады добавляются отдельно. Нельзя подставить сюда амплитудный коэффициент усиления без соответствующего множителя два. Фазовое условие задаёт продольные частоты, приблизительно $\nu_q\simeq qc/(2nL)$; межмодовый интервал с учётом дисперсии $\Delta\nu\simeq c/(2n_gL)$ для однородного заполнения. Поперечные профили, дифракция и перекрытие с накачкой дополнительно различают моды.

\paragraph{Насыщение одной моды.}
В простом стационарном приближении $g(I)=g_0/(1+I/I_{\mathrm{sat}})$. Генерация возможна при $g_0>g_{\mathrm{th}}$, а баланс даёт
\begin{equation}
 I_*=I_{\mathrm{sat}}(g_0/g_{\mathrm{th}}-1).
\end{equation}
Это локальная модель интенсивности при заданной эффективной накачке, а не универсальная формула выходной мощности любого лазера. Для динамики поля и инверсии полезна модель
\begin{equation}\label{eq:laserRate}
 \dot S=(GN-\tau_p^{-1})S,\qquad
 \dot N=R-N/\tau_N-GNS.
\end{equation}
Здесь $S$ и $N$ --- согласованно нормированные плотности фотонов и возбуждений, $G$ включает эффективное перекрытие, $\tau_p$ --- время жизни \emph{энергии} поля, $\tau_N$ --- релаксация возбуждений. Спонтанный источник опущен для анализа стационарной генерации; физически он запускает поле. Модель соответствует эффективной четырёхуровневой схеме с однородными распределениями и не заменяет общие уравнения Максвелла--Блоха.

Нулевой режим имеет $N=R\tau_N$; порог $R_{\mathrm{th}}=1/(G\tau_p\tau_N)$. Выше порога
\begin{equation}
 N_*=(G\tau_p)^{-1},\qquad S_*=\tau_p(R-R_{\mathrm{th}}).
\end{equation}
Инверсия фиксируется около порога: дополнительная накачка преимущественно переходит в фотоны. Линеаризация даёт
\begin{equation}
 J=\begin{pmatrix}0&GS_*\\-\tau_p^{-1}&-\tau_N^{-1}-GS_*\end{pmatrix},
 \quad \operatorname{tr}J<0,\quad\det J=GS_*/\tau_p>0.
\end{equation}
Режим устойчив в этой модели; при отрицательном дискриминанте возникают затухающие релаксационные колебания. При $\tau_p\ll\tau_N$ их частота порядка $\sqrt{(r-1)/(\tau_p\tau_N)}$, $r=R/R_{\mathrm{th}}$, если затухание мало по сравнению с этой частотой. Амплитуда пассивного поля убывает с декрементом $1/(2\tau_p)$, поэтому $Q=\omega\tau_p$.
